% ============================================================
%  backup.tex  —  FRAGMENT, included via \input from main.tex
%  Backup slides: MDP components that are UNCHANGED from v1 to v2.
%  Placed behind \appendix so they are excluded from the main
%  frame count and the outline. Uses \notebar (defined in abel.tex)
%  and colors/styles from main.tex.
% ============================================================

\appendix

% ============================================================
\section{Backup Slides}
% ============================================================

% ---- Backup: Action (UNCHANGED) ----------------------------
\begin{frame}{Action $\;a_t$ \hfill \small\textcolor{SkyBlue!70!black}{(unchanged)}}
  \small
  A single \textbf{discrete} choice over the visible queue:
  $$ \mathcal{A} = \{0, 1, 2, \dots, K\} = \text{Discrete}(K{+}1) $$
  \begin{itemize}
    \item $a_t = 0$ (\textbf{no-op}): leave the queue ordering unchanged.
    \item $a_t = k \in \{1,\dots,K\}$: \textbf{promote} the $k$-th visible queued
          job to the front, so it is considered first on the next dispatch.
  \end{itemize}
  \begin{block}{Key idea}
    The agent never dispatches jobs directly; it only \emph{reprioritises}.
    The underlying scheduler then greedily packs jobs that fit onto free nodes.
  \end{block}
  \notebar{$K$ = number of visible queued jobs \;\textbar\; $a_t$ = action at step $t$ \;\textbar\; $k$ = chosen queue position}
\end{frame}

% ---- Backup: Transition (UNCHANGED) ------------------------
\begin{frame}{Transition $\;P(s_{t+1}\mid s_t, a_t)$ \hfill \small\textcolor{SkyBlue!70!black}{(unchanged)}}
  \small
  Each step advances the simulation by $\Delta t$ seconds:
  \begin{enumerate}
    \item \textbf{Apply the action}: if $a_t = k > 0$, move that job to the head of the queue.
    \item \textbf{Step the scheduler} by $\Delta t$: dispatch queued jobs onto free
          nodes (greedy packing), advance running jobs, release nodes from jobs that finish.
    \item \textbf{New arrivals}: jobs arrive stochastically (Poisson-like
          inter-arrival times set by the \emph{offered load}).
  \end{enumerate}
  \begin{block}{Source of stochasticity}
    Job arrivals and the workload sampled each episode make the next state only
    \emph{partially} predictable from $(s_t, a_t)$.
  \end{block}
  \notebar{$\Delta t$ = simulation step (seconds) \;\textbar\; $a_t$ = action at step $t$ \;\textbar\; $k$ = promoted job index}
\end{frame}

% ---- Backup: Horizon & Objective (UNCHANGED) ---------------
\begin{frame}{Horizon \& Objective \hfill \small\textcolor{SkyBlue!70!black}{(unchanged)}}
  \small
  An episode \textbf{truncates} when either
  $$ t \geq T \quad(\text{24\,h elapsed}) \qquad\text{or}\qquad \text{queue \& running sets are empty (drained).} $$
  There is no absorbing failure state, so \texttt{terminated} is always \texttt{False};
  episodes end only by truncation.

  \begin{block}{Learning objective}
    The agent learns a policy $\pi_\theta(a_t \mid s_t)$ maximising the expected
    discounted return:
    $$ \max_{\theta}\; \mathbb{E}_{\pi_\theta}\!\left[\sum_{t=0}^{H-1} \gamma^{\,t}\, r_t\right] $$
    i.e.\ \textbf{keep the cluster busy} while draining the queue.
  \end{block}
  \notebar{$t,T$ = elapsed / horizon time \;\textbar\; $H$ = steps per episode \;\textbar\; $\gamma$ = discount factor \;\textbar\; $r_t$ = reward at step $t$}
\end{frame}
